\lecture{9}
\title{Introducing the MLE}

\begin{document}

\subsection{Types of Parameters, Identifiability}

\textbf{Definition.} In a statistical model,
the parameters we care about are called \textbf{parameters of interest},
and those we don't are called \textbf{nuisance parameters}.
A parameter $\theta$ is \textbf{identifiable} if
distinct values of $\theta$ necessarily yield distinct probability distributions.

\textbf{Example.} In the model
$\{\mathcal{N}(\mu, \sigma^2) \mid \mu \in \mathbb{R}, \sigma \in \mathbb{R}^+\}$,
we might say $\mu$ is a parameter of interest and $\sigma$ is a nuisance parameter.

\textbf{Example.} The response to a drug dosage is modeled by
$\mathcal{N}\left(E_0 + \frac{D \times E_{\text{max}}}{D + ED_{50}},\ \sigma^2\right)$,
for some parameter $\theta = (\sigma^2, E_0, E_{\text{max}}, ED_{50})$.
In this case, $\theta$ is not identifiable.

\subsection{Defining the Maximum Likelihood Estimator (MLE)}

Consider some (i.i.d.) random variable samples $X_1, \dots, X_n \sim \mathbb{P}$
and a parameter $\theta$.  Furthermore, some new notation:

\textbf{Notation.} Given a parameter $\theta$,
let its true, population value be denoted by $\theta^*$.

If the parameter $\theta$ can be expressed using $\mathbb{E}[\dots]$ only,
it's not hard to define an okay estimator.

\[ \text{(ex:)} ~~~~~~ \theta := \frac{\mu}{\sigma} =
\frac{\mathbb{E}[X_i]}{(\mathbb{E}[X_i^2] - \mathbb{E}[X_i]^2)^{1/2}} ~ \implies ~ \hat{\theta}_n :=
\frac{\frac{1}{n}\sum_{i = 1}^n X_i}{\left(
\frac{1}{n}\sum_{i = 1}^n X_i^2 - \left(\frac{1}{n}\sum_{i = 1}^n X_i\right)^2
\right)^{1/2}}. \]

More smartly, though, we should think to use the \textbf{maximum likelihood estimator (MLE)}.

\textbf{Definition.} Let $f_{\theta}(x)$ denote the PDF corresponding to
the distribution $\mathbb{P}_{\theta}$, and
consider some samples $X_1, \dots, X_n \sim \mathbb{P}_{\theta^*}$.
\begin{itemize}
    \item The \textbf{likelihood} function is defined by
    $L_n(\theta) := \prod_{i = 1}^n f_{\theta}(X_i)$.
    \item The \textbf{log likelihood} function is just $\ell_n(\theta) := \log L_n(\theta)
    = \sum_{i = 1}^n \log f_{\theta}(X_i)$.
    \item The \textbf{maximum likelihood estimator (MLE)} $\hat{\theta}_n$ is the value
    of $\theta$ for which $L_n(\theta)$ (or $\ell_n(\theta)$) is maximized.
\end{itemize}
In other words, the MLE $\hat{\theta}_n$
is the value for the parameter $\theta$ whose distribution $\mathbb{P}_{\hat{\theta}_n}$
is most ``compatible'' with the samples taken from $\mathbb{P}_{\theta^*}$,
which would suggest that $\mathbb{P}_{\hat{\theta}_n}$
might be a good representation of $\mathbb{P}_{\theta^*}$.

\begin{remark}
    Reason through the following and check that they make sense.
    \begin{itemize}
        \item For the model $\{\mathrm{Ber}(\theta) \mid \theta \in [0, 1]\}$,
        the MLE is $\hat{\theta} := \frac{1}{n} \sum_{i = 1}^n X_i$.
        \item For the model $\{\mathcal{N}(\theta, 1) \mid \theta \in \mathbb{R}\}$,
        the MLE is $\hat{\theta} := \frac{1}{n} \sum_{i = 1}^n X_i$.
        \item For the model $\{\mathrm{Unif}([0, \theta]) \mid \theta \geq 0\}$,
        the MLE is $\hat{\theta} := \max X_i$.
    \end{itemize}
\end{remark}

Note that $L_n(\theta)$ and $\ell_n(\theta)$ are
\emph{deterministic} functions whose behaviors are determined by \emph{random} samples
from $\mathbb{P}_{\theta^*}$.
If one seeks a likelihood function that is instead determined by $\theta^*$ itself,
one may opt for\dots

\textbf{Definition.} The \textbf{population log likelihood} function is 
 $\ell(\theta) := \mathbb{E}_{\theta^*}[\log f_{\theta}(X)]$.
(The $\mathbb{E}_{\theta^*}$ notation
indicates that $X \sim \mathbb{P}_{\theta^*}$.)

\end{document}